% !!! GENERATED FILE — DO NOT EDIT !!!
% Produced by manuscript/generate_catalog_table.jl from, in order of authority:
%   packages/VBFCatalog/data-raw/apn_catalog/n{4,5,6,7,8,9,10,12}_apn.toml
%   packages/VBFCatalog/data-raw/bibliography.bib
%   packages/VBFCatalog/BFA_catalog.sqlite   (functions, bulk_corpora,
%                                             equivalence_classes, class_members)
%   packages/VBFCatalog/corpora/*.sqlite
%   packages/VBFCatalog/data-build/*.jl, packages/VBF/src/**.jl
%   papers/<year>/<Paper_ID>/README.md
% Regenerate with:  julia --project=. -t auto manuscript/generate_catalog_table.jl
\item[a] \texttt{n6\_Kasami\_x48} is not a distinct CCZ class: $48 \equiv 3\cdot 2^4 \bmod 63$, and the catalog records \texttt{are\_ea\_equivalent} and \texttt{are\_ccz\_equivalent} both returning true against \texttt{n6\_Gold\_x3}. The in-repo CCZ index places both, and four further $n=6$ entries, in one class.
\item[b] The two $n=9$ entries sit in distinct in-repo CCZ classes, and in classes distinct from every $n=9$ Gold entry. Inequivalence to the non-quadratic AB power maps at $n=9$ is taken from the source paper's $\Gamma$-rank argument; the catalog states that an exhaustive in-repo CCZ decision is intractable at $n=9$.
\item[c] The reproduction review of this paper records that its own test suite could not be executed for want of a Julia toolchain, so the novelty comparison it reports is not an in-repo recomputation.
\item[d] The reproduction review reports \texttt{classify\_ccz} returning \texttt{:known} for all three $n=8$ representatives, so no novelty is asserted for them here.
\item[e] CCZ-inequivalence of the Kim function to both $n=6$ power entries is recorded as verified in-repo. The Dillon permutation itself has algebraic degree 4 and therefore lies outside the quadratic-only in-repo CCZ decision.
\item[f] Inequivalence to Gold, hexanomial, Taniguchi and Zhou--Pott at $n=8$ was confirmed independently; the paper's further claim of inequivalence to two external $n=8$ tables is recorded as unverifiable here because those corpora are not in the repository.
\item[g] Novelty is established against the eleven known-family rows of the paper's own Table 3, all of which were recomputed. The check is not exercised by the reproduction's test suite.
\item[h] The published $\Gamma$-ranks are refused by the library's $n\ge 8$ design-rank guard. The repository is inconsistent on this point: the reproduction review first records the value as verified, then reverses that in a later re-review, while the catalog entry states it is not verified. The guard in \texttt{packages/VBF} is unchanged.
\item[i] This entry is EA-equivalent, hence CCZ-equivalent, to Budaghyan--Carlet--Leander 2009 eq.\,(6) with $a=\alpha$ over $\mathrm{GF}(2^8)$. The generator decides this on every run with \texttt{are\_ea\_equivalent}, which is complete on the quadratic APN class; the two lookup tables are not identical, because the constructions use different moduli. The reproduction review of the 2009 paper had established only that the two ortho-derivative mugshots agree --- a necessary condition --- and had left the decision open. \emph{The divergence is against this catalog's provenance record}, which attributes the entry to 2025 alone and carries no prior-art key. It is not a divergence against the 2025 paper, which claims only EA-inequivalence to $x^3$ and asserts no novelty; that is why the cell carries no E-code. The catalog entry needs a cross-reference to the 2009 construction. This table reports the collision and does not change the catalog.
\item[j] The entry itself states that absence from this catalog is not a claim of novelty in the literature, and that the paper claims only EA-inequivalence to $x^3$.
\item[k] The corpus registry labels this set ``8{,}157 APN functions'' while the side store holds \corpusYWL{} members after one non-quadratic member was excluded at import. The repository carries no analysis that reconciles the two figures.
\item[l] Every member is indexed by its ortho-derivative signature, which this repository computes from the function. That signature is a necessary condition for CCZ-equivalence, not a decision, and the members are not partitioned into confirmed CCZ classes.
\item[m] The count of pairwise distinct EA classes reproduces exactly; that the classes are new relative to the earlier corpora is not recomputed here.
\item[n] Four members were shown pairwise CCZ-inequivalent by their ortho-derivative differential spectra. The remaining members carry ortho-derivative signatures only.
\item[o] The reproduction records the switching sub-family as not separable from the corpus without the authors' provenance labels, and the paper's $n=9$ and $n=10$ instances as absent from the repository.
\item[p] These sets carry no \texttt{bib\_key} in the catalog: \texttt{bulk\_corpora.bib\_key} and \texttt{functions.bib\_key} are \textsc{null} for every one of their rows, so no citation is printed rather than one invented. The upstream sources are named in \texttt{data-build/sources.toml} as the Simon--Berg thesis and the \texttt{sboxU} repository.
\item[q] The 27 class representatives are named entries with an in-repo confirmed CCZ class. The assignment of all \corpusSB{} members to those 27 classes is taken from the source and is not recomputed here.
\item[r] The reproduction review of this source records that its test suite could not be executed for want of a Julia toolchain, so no figure it reports is an in-repo recomputation.
\item[s] Two of the imported representatives are not quadratic and therefore lie outside the in-repo CCZ decision.
\item[t] The reproduction reports rows 4 and 10 of the source table as the same polynomial, so the table lists 14 distinct entries where the classification has 15 classes; the missing class was recovered by exhaustive search in the reproduction.
